
%+TITLE: Repaso de álgebra lineal
%+DESCRIPTION: Repaso de los resultados del álgebra lineal que necesitaremos
\documentclass[11pt,a4paper]{article}

\usepackage[utf8]{inputenc}
\usepackage[T1]{fontenc}
\usepackage[spanish,es-nodecimaldot]{babel}
\usepackage{lmodern}
\usepackage{amsmath,amssymb,mathtools}
\usepackage{graphicx}
\usepackage{xcolor}
\usepackage{geometry}
\usepackage{enumitem}
\usepackage{microtype}

\geometry{margin=2.3cm}
\setlength{\parindent}{0pt}
\setlength{\parskip}{0.65em}
\setlist[itemize]{leftmargin=1.6em,itemsep=0.25em,topsep=0.25em}

\definecolor{title-red}{RGB}{165,25,20}
\definecolor{concept-green}{RGB}{20,125,45}
\newcommand{\concept}[1]{\textcolor{concept-green}{#1}}

\newcommand{\ket}[1]{\left|#1\right\rangle}
\newcommand{\bra}[1]{\left\langle#1\right|}
\newcommand{\braket}[2]{\left\langle #1\middle|#2\right\rangle}
\newcommand{\hc}{\mathrm{h.c.}}
\newcommand{\Tr}{\operatorname{Tr}}
\newcommand{\Diffs}{\mathrm{Diffs}}

\begin{document}

Para entender las simetrías necesitamos la teoría de grupos y el álgebra lineal. Los grupos los estudiaremos en su debido momento. Ahora repasamos el álgebra lineal.

Del estudio de fuerzas y posiciones en mecánica tenemos una idea de vectores. Lo importante de los vectores es que se suman de forma lineal (se ``superponen''):
\[
  3(4\hat x+\hat y)+(2\hat x)=14\hat x+3\hat y
  \qquad\longrightarrow\qquad \text{lineal}.
\]
Si en cambio fuera
\[
  3(4\hat x+\hat y)+(2\hat x)=16\hat x+3\hat y,
\]
no sería lineal.

Un \concept{espacio vectorial} es un conjunto de elementos, $\{\ket{\chi_1},\ket{\chi_2},\ldots\}$, junto con números complejos y operaciones de suma y multiplicación por un complejo tales que

\begin{itemize}
  \item $\ket{\chi_1}+\ket{\chi_2}$ y $\alpha\ket{\chi}$ son vectores.

  \item Hay un vector cero $\ket{0}$ tal que $\ket{\chi}+\ket{0}=\ket{\chi}$, y cada vector tiene un inverso $-\ket{\chi}$ tal que $\ket{\chi}-\ket{\chi}=\ket{0}$. Además, $0\ket{\chi}=\ket{0}$.

  \item Todas las operaciones son asociativas:
  \[
    \ket{\chi_1}+(\ket{\chi_2}+\ket{\chi_3})
    =(\ket{\chi_1}+\ket{\chi_2})+\ket{\chi_3},
    \qquad
    \alpha(\beta\ket{\chi})=(\alpha\beta)\ket{\chi}.
  \]

  \item
  \[
    \alpha_1\ket{\chi}+\alpha_2\ket{\chi}
      =(\alpha_1+\alpha_2)\ket{\chi},
    \qquad
    \alpha\ket{\chi_1}+\alpha\ket{\chi_2}
      =\alpha(\ket{\chi_1}+\ket{\chi_2}),
  \]
  lo que expresa la \concept{linealidad}.
\end{itemize}

$\alpha_1\ket{\chi_1}+\alpha_2\ket{\chi_2}$ se llama una \concept{combinación lineal}. Un conjunto de vectores $\{\ket{\chi_1},\ldots,\ket{\chi_n}\}$ es \concept{linealmente independiente} si
\[
  \alpha_1\ket{\chi_1}+\cdots+\alpha_n\ket{\chi_n}=\ket{0}
\]
implica $\alpha_1=\cdots=\alpha_n=0$. Es decir, ningún vector del conjunto se puede escribir como una combinación lineal de los demás.

Por ejemplo, $\{\ket{\chi},\ket{y},2\ket{\chi}+\ket{y}\}$ no es linealmente independiente:
\[
  (2\ket{\chi}+\ket{y})-2\ket{\chi}-\ket{y}=\ket{0}.
\]

Dado un conjunto de vectores $S=\{\ket{\chi_1},\ldots,\ket{\chi_n}\}$, el espacio generado por $S$ es el espacio vectorial de todas las posibles combinaciones lineales de elementos de $S$. Una \concept{base} es un conjunto de vectores linealmente independientes que genera todo el espacio vectorial.

Vamos a considerar espacios vectoriales con un \concept{producto interno}. Es decir, una operación que toma dos elementos del espacio vectorial y da un número tal que
\begin{align*}
  (\ket v,\ket w)&=(\ket w,\ket v)^*,\\
  (\ket v,\alpha\ket{w_1}+\beta\ket{w_2})
    &=\alpha(\ket v,\ket{w_1})+\beta(\ket v,\ket{w_2}),\\
  (\alpha\ket{v_1}+\beta\ket{v_2},\ket w)
    &=\alpha^*(\ket{v_1},\ket w)+\beta^*(\ket{v_2},\ket w).
\end{align*}

Para hacer la notación más sencilla vamos a escribir $(\ket v)^*=\bra v$, tal que
\[
  (\alpha\ket v+\beta\ket w)^*=\alpha^*\bra v+\beta^*\bra w.
\]
De esta manera, estas reglas son fáciles de recordar si definimos $(\ket a,\ket b)=\braket a b$. Por ejemplo, el producto interno entre $\alpha\ket{v_1}+\beta\ket{v_2}$ y $\ket w$ es
\[
  (\alpha^*\bra{v_1}+\beta^*\bra{v_2})\ket w
  =\alpha^*\braket{v_1}{w}+\beta^*\braket{v_2}{w}.
\]

Una base $\{\ket{e_1},\ldots,\ket{e_n}\}$ es ortonormal si
\[
  \braket{e_i}{e_j}=\delta_{ij},
  \qquad
  \delta_{ij}=\begin{cases}
    1,&i=j,\\
    0,&i\neq j,
  \end{cases}
\]
donde $\delta_{ij}$ es la delta de Kronecker.

Un \concept{operador} o \concept{transformación lineal} $A$ es una función que toma un elemento del espacio vectorial y da otro elemento del espacio vectorial tal que
\[
  A(\alpha\ket{v_1}+\beta\ket{v_2})
  =\alpha A\ket{v_1}+\beta A\ket{v_2}.
\]

Notemos que definimos $A(\ket v)=A\ket v$ y lo vemos como ``$A$ actuando sobre $\ket v$''.

Con una base, podemos escribir las \concept{componentes} del vector:
\[
  \ket v=v_1\ket{e_1}+\cdots+v_n\ket{e_n},
  \qquad\text{o también}\qquad
  \ket v=\begin{pmatrix}v_1\\ \vdots\\ v_n\end{pmatrix}.
\]
Si la base es ortonormal, $\braket{e_1}{v}=v_1$, o en general $v_i=\braket{e_i}{v}$.

\[
  A\ket v=A(v_1\ket{e_1}+\cdots+v_n\ket{e_n})
  =v_1A\ket{e_1}+\cdots+v_nA\ket{e_n}.
\]
Por tanto, si sabemos cómo actúa $A$ sobre los elementos de la base, podemos calcular $A\ket v$.

Las componentes del nuevo vector $A\ket{e_j}$ son $\bra{e_i}A\ket{e_j}$, que llamamos $A_{ij}$. De esta manera, $A$ se puede escribir como una matriz.

Una notación útil es
\[
  \ket v=v_1\ket{e_1}+\cdots+v_n\ket{e_n}
  =\sum_{i=1}^n v_i\ket{e_i}\equiv v_i\ket{e_i}.
\]
En la \concept{notación de Einstein}, los índices repetidos se suman.

Calcular las componentes de $A\ket v$ es fácil:
\begin{align*}
  \bra{e_i}A\ket v
  &=\bra{e_i}A\bigl(\ket{e_j}\braket{e_j}{v}\bigr)\\
  &=\bra{e_i}A\ket{e_j}\braket{e_j}{v}
   =A_{ij}v_j
   =
   \begin{pmatrix}
     A_{11}&\cdots&A_{1n}\\
     \vdots&\ddots&\vdots\\
     A_{n1}&\cdots&A_{nn}
   \end{pmatrix}
   \begin{pmatrix}v_1\\ \vdots\\ v_n\end{pmatrix}.
\end{align*}
Para obtener las nuevas componentes usamos la clásica multiplicación matriz por vector.

Definamos algunas operaciones útiles para matrices.

\concept{Traza}:
\[
  \Tr A=A_{ii},
  \qquad
  \Tr(ABC)=A_{ij}B_{jk}C_{ki}=C_{ki}A_{ij}B_{jk}=\Tr(CAB).
\]

\concept{Determinante}:
\[
  \det A=\sum_{\sigma}\operatorname{sign}(\sigma)
  A_{1\sigma(1)}\cdots A_{n\sigma(n)},
\]
donde la suma es sobre las permutaciones de $1,\ldots,n$.

Para aprender a calcular los determinantes refiero a los libros de álgebra lineal.

Se puede demostrar que
\[
  \det(AB)=\det A\,\det B.
\]

Definimos la \concept{traspuesta} mediante
\[
  (A^T)_{ij}=A_{ji};
\]
resulta que $\det A^T=\det A$.

El operador \concept{inverso} $A^{-1}$ es tal que
\[
  A^{-1}A=\mathbb{1}.
\]
Esto quiere decir que $\det A^{-1}=1/\det A$.

Definimos
\[
  (A\ket v)^*=\bra v A^\dagger,
\]
que introduce el hermítico conjugado. En componentes,
\[
  (A^\dagger)_{ij}
  =\bra{e_i}A^\dagger\ket{e_j}
  =\bigl(\bra{e_j}A\ket{e_i}\bigr)^*
  =A_{ji}^*,
\]
es decir, traspuesto y complejo conjugado.

Un operador es \concept{hermítico} si $A^\dagger=A$.

Un operador es \concept{unitario} si $A^\dagger=A^{-1}$. En los reales, un operador es \concept{ortogonal} si $A^T=A^{-1}$.

Un espacio vectorial puede tener muchas bases diferentes. Los elementos de una base $\{\ket{e_i'}\}$ se pueden expresar en términos de otra $\{\ket{e_i}\}$:
\[
  \ket{e_i'}=\ket{e_j}\braket{e_j}{e_i'},
  \qquad
  S_{ji}=\braket{e_j}{e_i'},
\]
donde $S$ es la \concept{matriz de cambio de base}.

Si la base $\{\ket{e_i'}\}$ es ortonormal,
\[
  \delta_{ij}
  =\braket{e_i'}{e_j'}
  =\braket{e_i'}{e_k}\braket{e_k}{e_j'}
  =S^\dagger_{ik}S_{kj},
  \qquad\Longrightarrow\qquad
  S^\dagger S=\mathbb{1}.
\]

También podemos expresar $A$ en la base $\{\ket{e_i'}\}$:
\begin{align*}
  A'_{ij}
  &=\bra{e_i'}A\ket{e_j'}\\
  &=\braket{e_i'}{e_k}\bra{e_k}A\ket{e_l}\braket{e_l}{e_j'}
   =S^\dagger_{ik}A_{kl}S_{lj},
\end{align*}
lo que es una \concept{transformación de similaridad}.

Ahora queremos saber si podemos encontrar una base en la cual $A$ tome una forma diagonal, es decir,
\[
  \bra{e_i'}A\ket{e_j'}=D_{ij},
  \qquad
  D=\begin{pmatrix}
    \lambda_1&&0\\
    &\ddots&\\
    0&&\lambda_n
  \end{pmatrix}.
\]

Trabajar en esta base simplifica las cosas. ¿Cuándo podemos hacer algo así?

Encontrar este cambio de base se llama \concept{diagonalizar} el operador. El hecho de que el operador sea diagonal quiere decir
\[
  A\ket{e_i'}=\lambda_{(i)}\ket{e_i'}
\]
(los índices entre paréntesis no se suman). Decimos que $\ket{e_i'}$ es un \concept{autovector} con \concept{autovalor} $\lambda_{(i)}$.

Para encontrar los autovalores y autovectores necesitamos resolver
\begin{align*}
  A\ket v=\lambda\ket v
  &\Longrightarrow \bra{e_i}A\ket v=\lambda\braket{e_i}{v}\\
  &\Longrightarrow \bra{e_i}A\ket{e_j}\braket{e_j}{v}
     =\lambda\braket{e_i}{v}\\
  &\Longrightarrow A_{ij}V_j=\lambda V_i\\
  &\Longrightarrow (A_{ij}-\lambda\delta_{ij})V_j=0.
\end{align*}
Son $n$ ecuaciones con $n$ componentes de $V$ si fijamos $\lambda$.

Claramente, $V=(0,\ldots,0)^T$ es una solución. No queremos que esta sea la única solución. Es decir, queremos que el sistema de ecuaciones sea degenerado.

Un sistema de ecuaciones $B_{ij}X_j=0$ es degenerado cuando una de las ecuaciones, digamos la primera, $B_{1j}X_j=0$, se puede escribir como una combinación lineal de las otras. En otras palabras, esa ecuación no es independiente y no hay suficientes ecuaciones para las variables.

Escribamos explícitamente la primera fila como combinación lineal de las demás:
\begin{align}
B &=
\begin{pmatrix}
 \alpha_2b_{21}+\cdots+\alpha_nb_{n1} & \cdots &
 \alpha_2b_{2n}+\cdots+\alpha_nb_{nn}\\
 b_{21}&\cdots&b_{2n}\\
 \vdots&\ddots&\vdots\\
 b_{n1}&\cdots&b_{nn}
\end{pmatrix} \\
&=\alpha_2
\underbrace{\begin{pmatrix}
 b_{21}&\cdots&b_{2n}\\
 b_{21}&\cdots&b_{2n}\\
 \vdots&\ddots&\vdots\\
 b_{n1}&\cdots&b_{nn}
\end{pmatrix}}_{B^{(2)}}
+\cdots+
\alpha_n
\underbrace{\begin{pmatrix}
 b_{n1}&\cdots&b_{nn}\\
 b_{21}&\cdots&b_{2n}\\
 \vdots&\ddots&\vdots\\
 b_{n1}&\cdots&b_{nn}
\end{pmatrix}}_{B^{(n)}}
\end{align}
Por linealidad del determinante en una fila,
\[
  \det B=\alpha_2\det B^{(2)}+\cdots+\alpha_n\det B^{(n)}.
\]

Pero el determinante de una matriz con filas repetidas es cero. (El determinante es la suma de todas las permutaciones con un signo; si hay repetición da cero.) Entonces, el sistema $B_{ij}X_j=0$ es degenerado si y sólo si $\det B=0$.

Por tanto,
\[
  (A-\lambda\mathbb{1})V=0
\]
tiene soluciones diferentes de $V=(0,\ldots,0)^T$ si y sólo si
\[
  \det(A-\lambda\mathbb{1})=0,
\]
la \concept{ecuación característica}.

Esta ecuación característica, escrita explícitamente, es una ecuación de orden $n$ para $\lambda$. Esta tendrá $n$ soluciones para $\lambda$. Cada una se pone en $Av=\lambda v$ para encontrar $n$ vectores (uno por $\lambda$).
\end{document}
